\section*{Sets and Indices}

\begin{itemize}
    \item $F$: set of fruit types, indexed by $f$, with
    \[
    F=\{\text{apples},\text{pears},\text{oranges},\text{lemons}\}.
    \]
    \item $S$: set of growing seasons, indexed by $s$, with
    \[
    S=\{\text{spring},\text{autumn}\}.
    \]
\end{itemize}

\section*{Parameters}

\begin{itemize}
    \item $p_f$: profit per acre for fruit $f$ (code: \texttt{profit[f]}), from
    \texttt{profit\_per\_acre\_apples}, \texttt{profit\_per\_acre\_pears},
    \texttt{profit\_per\_acre\_oranges}, \texttt{profit\_per\_acre\_lemons}.
    \item $A$: total farm area available in each season (code: \texttt{total\_area}), from \texttt{total\_farm\_area}.
    \item $\rho^{ap}$: minimum apples-to-pears annual area ratio (code: \texttt{min\_a\_to\_p}), from \texttt{min\_apples\_to\_pears\_ratio}.
    \item $\rho^{al}$: minimum apples-to-lemons annual area ratio (code: \texttt{min\_a\_to\_l}), from \texttt{min\_apples\_to\_lemons\_ratio}.
    \item $\rho^{ol}$: oranges-to-lemons annual area ratio (code: \texttt{o\_to\_l}), from \texttt{oranges\_to\_lemons\_ratio}.
    \item $K$: maximum number of fruit types allowed annually (code: \texttt{max\_types}), from \texttt{max\_fruit\_types}.
    \item $w_{fs}$: water required per acre of fruit $f$ in season $s$ (code: \texttt{water\_per\_acre[(f,s)]}), from the corresponding \texttt{water\_per\_acre\_<fruit>\_<season>} entries.
    \item $W^{sp}$: spring water cap (code: \texttt{water\_cap\_spring}), from \texttt{water\_cap\_spring}.
    \item $W^{au}$: autumn water cap (code: \texttt{water\_cap\_autumn}), from \texttt{water\_cap\_autumn}.
    \item $W^{tot}$: annual total water budget (code: \texttt{total\_water\_budget}), from \texttt{total\_water\_budget}.
    \item $\underline{W}$: minimum annual water required for any activated fruit type (code: \texttt{min\_annual\_water\_per\_fruit}), from \texttt{min\_annual\_water\_per\_fruit}.
    \item $R$: diversification reward (code: \texttt{diversification\_reward}), from \texttt{diversification\_reward}.
    \item $M^A$: area big-$M$ constant (code: \texttt{big\_M\_area}),
    \[
    M^A = A\cdot |S|.
    \]
    \item $M^W$: water big-$M$ constant (code: \texttt{big\_M\_water}),
    \[
    M^W = \Big(\max_{f\in F,\ s\in S} w_{fs}\Big)\, M^A.
    \]
\end{itemize}

\section*{Decision Variables}

\begin{itemize}
    \item $a_{fs}\ge 0$: acres of fruit $f$ grown in season $s$ (code: \texttt{area[(f,s)]}).
    \item $u_{fs}\ge 0$: water used by fruit $f$ in season $s$ (code: \texttt{water[(f,s)]}).
    \item $y_f\in\{0,1\}$: 1 if fruit $f$ is counted as an active fruit type in the annual plan (code: \texttt{y[f]}).
    \item $z_f\in\{0,1\}$: binary indicator linked to annual water use of fruit $f$ (code: \texttt{z[f]}).
    \item $v\in\{0,1\}$: diversification-reward indicator (code: \texttt{v}).
\end{itemize}

For convenience, define annual totals
\[
A_f := \sum_{s\in S} a_{fs}, \qquad U_f := \sum_{s\in S} u_{fs}.
\]

\section*{Objective}

\[
\max \sum_{f\in F}\sum_{s\in S} p_f\, a_{fs} + R\, v.
\]

\section*{Constraints}

Seasonal land availability:
\[
\sum_{f\in F} a_{fs} \le A \qquad \forall s\in S.
\]

Annual crop-mix rules:
\[
A_{\text{apples}} \ge \rho^{ap}\, A_{\text{pears}},
\]
\[
A_{\text{apples}} \ge \rho^{al}\, A_{\text{lemons}},
\]
\[
A_{\text{oranges}} = \rho^{ol}\, A_{\text{lemons}}.
\]

Activation of fruit types through annual area:
\[
A_f \le M^A\, y_f \qquad \forall f\in F.
\]

Maximum number of fruit types:
\[
\sum_{f\in F} y_f \le K.
\]

Seasonal water-definition constraints:
\[
u_{fs} = w_{fs}\, a_{fs} \qquad \forall f\in F,\ \forall s\in S.
\]

Seasonal water caps:
\[
\sum_{f\in F} u_{f,\text{spring}} \le W^{sp},
\]
\[
\sum_{f\in F} u_{f,\text{autumn}} \le W^{au}.
\]

Annual water budget:
\[
\sum_{f\in F}\sum_{s\in S} u_{fs} \le W^{tot}.
\]

Water-based activation and minimum operating threshold:
\[
U_f \le M^W\, z_f \qquad \forall f\in F,
\]
\[
U_f \ge \underline{W}\, z_f \qquad \forall f\in F.
\]

Consistency between area-based and water-based activation:
\[
z_f = y_f \qquad \forall f\in F.
\]

Diversification reward logic. Let
\[
N := \sum_{f\in F} y_f.
\]
Then
\[
2v \le N,
\]
\[
N \le 1 + 4v.
\]

\section*{Notes on Implementation}

\begin{itemize}
    \item The model is implemented as a mixed-integer linear program in \texttt{current\_heuristic.py} and solved with \texttt{pulp.GUROBI\_CMD}.
    \item Profit is earned per acre and is not season-dependent in the code; only water coefficients vary by season.
    \item The farm area constraint is enforced separately in spring and autumn, so the same physical land may be reused across seasons.
    \item The annual ratio constraints are applied to total annual acreage $A_f=\sum_s a_{fs}$, not season by season.
    \item A fruit can have positive area only if $y_f=1$ because of $A_f\le M^A y_f$. Since $z_f=y_f$ and $U_f\ge \underline{W} z_f$, any activated fruit must also consume at least \texttt{min\_annual\_water\_per\_fruit} over the year.
    \item The diversification logic is encoded exactly as in the source:
    \[
    2v \le \sum_f y_f,\qquad \sum_f y_f \le 1+4v.
    \]
    With four fruit candidates and binary $v$, this makes $v=1$ when at least two fruit types are active, and forces $v=0$ when at most one fruit type is active.
    \item The code computes
    \[
    M^A=A|S|,\qquad
    M^W=\Big(\max_{f,s} w_{fs}\Big)M^A,
    \]
    and these exact big-$M$ values should be retained to match the implementation.
\end{itemize}
